\documentclass[aspectratio=169,11pt]{beamer}

\usetheme{Madrid}
\usecolortheme{seahorse}
\setbeamertemplate{navigation symbols}{}
\setbeamertemplate{footline}[frame number]

\usepackage[utf8]{inputenc}
\usepackage[T5]{fontenc}
\usepackage[vietnamese,english]{babel}
\usepackage{array}
\usepackage{booktabs}
\usepackage{xcolor}
\usepackage{listings}
\usepackage{tikz}
\usetikzlibrary{arrows.meta,positioning,shapes.geometric}

\definecolor{Primary}{HTML}{0F766E}
\definecolor{Accent}{HTML}{2563EB}
\definecolor{Warn}{HTML}{D97706}
\definecolor{Dark}{HTML}{0F172A}
\definecolor{Soft}{HTML}{F1F5F9}

\setbeamercolor{title}{fg=white,bg=Primary}
\setbeamercolor{frametitle}{fg=white,bg=Primary}
\setbeamercolor{structure}{fg=Accent}
\setbeamercolor{block title}{fg=white,bg=Accent}
\setbeamercolor{block body}{fg=Dark,bg=Soft}
\setbeamercolor{alerted text}{fg=Warn}

\lstdefinestyle{code}{
  basicstyle=\ttfamily\scriptsize,
  keywordstyle=\color{Accent},
  commentstyle=\color{Primary},
  stringstyle=\color{Warn},
  breaklines=true,
  frame=single,
  rulecolor=\color{Accent!35},
  backgroundcolor=\color{Soft}
}
\lstset{style=code}

\title[Academic Appointment]{Academic Appointment System}
\subtitle{Hệ thống đặt lịch tư vấn học thuật giữa sinh viên và giảng viên}
\author{Nguyễn Văn Sinh}
\institute{MSSV: 2331200098 \quad Lớp: CNTT-K65}
\date{\today}

\begin{document}

\begin{frame}
  \titlepage
\end{frame}

\begin{frame}{Nội dung trình bày}
  \tableofcontents
\end{frame}

\section{Bài toán và giải pháp}

\begin{frame}{Bài toán}
  Trong môi trường học thuật, việc đặt lịch gặp giảng viên thường diễn ra qua tin nhắn, email hoặc trao đổi trực tiếp.

  \begin{block}{Các vấn đề chính}
    \begin{itemize}
      \item Sinh viên không biết chính xác khung giờ rảnh của giảng viên.
      \item Giảng viên khó kiểm soát nhiều yêu cầu đặt lịch cùng lúc.
      \item Dễ trùng lịch, quên phản hồi hoặc mất lịch sử xử lý.
      \item Admin thiếu công cụ tập trung để quản lý tài khoản và vai trò.
    \end{itemize}
  \end{block}
\end{frame}

\begin{frame}{Giải pháp đề xuất}
  \textbf{Academic Appointment System} chuẩn hóa quy trình đặt lịch học thuật thành một workflow có trạng thái và phân quyền rõ ràng.

  \begin{columns}[T]
    \begin{column}{0.5\textwidth}
      \begin{itemize}
        \item Giảng viên tạo slot tư vấn khả dụng.
        \item Sinh viên tìm giảng viên và đặt lịch.
        \item Giảng viên xác nhận, từ chối hoặc hủy lịch.
      \end{itemize}
    \end{column}
    \begin{column}{0.5\textwidth}
      \begin{itemize}
        \item Sinh viên theo dõi trạng thái và nhận thông báo.
        \item Admin quản lý user, role và trạng thái tài khoản.
        \item Hệ thống deploy được trên hạ tầng thật.
      \end{itemize}
    \end{column}
  \end{columns}

  \begin{alertblock}{Mục tiêu}
    Giảm thao tác thủ công, tránh trùng lịch và tập trung hóa quản lý appointment.
  \end{alertblock}
\end{frame}

\section{Chức năng và workflow}

\begin{frame}{Vai trò người dùng}
  \begin{columns}[T]
    \begin{column}{0.32\textwidth}
      \begin{block}{Student}
        \begin{itemize}
          \item Đăng ký, đăng nhập
          \item Xem giảng viên
          \item Xem slot còn trống
          \item Tạo và hủy lịch hẹn
          \item Nhận thông báo
        \end{itemize}
      \end{block}
    \end{column}
    \begin{column}{0.32\textwidth}
      \begin{block}{Lecturer}
        \begin{itemize}
          \item Quản lý hồ sơ
          \item Tạo slot tư vấn
          \item Xem request
          \item Confirm / Reject / Cancel
          \item Phản hồi cho sinh viên
        \end{itemize}
      \end{block}
    \end{column}
    \begin{column}{0.32\textwidth}
      \begin{block}{Admin}
        \begin{itemize}
          \item Xem danh sách user
          \item Xem chi tiết user
          \item Quản lý role
          \item Khóa / mở tài khoản
          \item Điều phối quyền truy cập
        \end{itemize}
      \end{block}
    \end{column}
  \end{columns}
\end{frame}

\begin{frame}{Chức năng chính}
  \begin{tabular}{p{0.28\textwidth} p{0.62\textwidth}}
    \toprule
    \textbf{Nhóm chức năng} & \textbf{Mô tả} \\
    \midrule
    Authentication & Đăng ký, đăng nhập, lấy thông tin user hiện tại \\
    Authorization & Phân quyền theo Admin, Student, Lecturer \\
    Appointment & Tạo, xem, cập nhật trạng thái, hủy lịch hẹn \\
    Availability Slot & Giảng viên tạo slot, sinh viên xem slot khả dụng \\
    Notification & Tạo thông báo khi lịch hẹn thay đổi \\
    Profile \& Avatar & Cập nhật hồ sơ và upload avatar qua Supabase \\
    Admin & Quản lý user, role và trạng thái hoạt động \\
    \bottomrule
  \end{tabular}
\end{frame}

\begin{frame}{Luồng đặt lịch}
  \centering
  \begin{tikzpicture}[
    node distance=1.15cm,
    box/.style={rectangle, rounded corners, draw=Accent, fill=Accent!8, text width=7.4cm, align=center, minimum height=0.75cm},
    arrow/.style={-Latex, thick, draw=Primary}
  ]
    \node[box] (a) {Lecturer tạo slot rảnh};
    \node[box, below=of a] (b) {Student xem lecturer detail};
    \node[box, below=of b] (c) {Student chọn slot và gửi yêu cầu};
    \node[box, below=of c] (d) {Appointment = \texttt{Pending}};
    \node[box, below=of d] (e) {Lecturer \texttt{Confirmed} / \texttt{Rejected} / \texttt{Cancelled}};
    \node[box, below=of e] (f) {Student nhận notification và theo dõi trạng thái};
    \draw[arrow] (a) -- (b);
    \draw[arrow] (b) -- (c);
    \draw[arrow] (c) -- (d);
    \draw[arrow] (d) -- (e);
    \draw[arrow] (e) -- (f);
  \end{tikzpicture}
\end{frame}

\begin{frame}{Trạng thái appointment}
  \begin{columns}[T]
    \begin{column}{0.5\textwidth}
      \begin{block}{Các trạng thái chính}
        \begin{itemize}
          \item \texttt{Pending}: chờ giảng viên xử lý.
          \item \texttt{Confirmed}: giảng viên đã xác nhận.
          \item \texttt{Rejected}: giảng viên từ chối.
          \item \texttt{Cancelled}: lịch bị hủy.
        \end{itemize}
      \end{block}
    \end{column}
    \begin{column}{0.5\textwidth}
      \begin{block}{Quy tắc nghiệp vụ}
        \begin{itemize}
          \item Student tạo appointment ở trạng thái \texttt{Pending}.
          \item Lecturer chỉ cập nhật lịch liên quan đến mình.
          \item Khi reject/cancel, slot được mở lại nếu phù hợp.
        \end{itemize}
      \end{block}
    \end{column}
  \end{columns}
\end{frame}

\begin{frame}[fragile]{Kiểm soát trùng lịch}
  Project xử lý trùng slot ở cả service layer và database constraint.

  \begin{block}{Service layer}
    \begin{itemize}
      \item Kiểm tra slot tồn tại và chưa bị xóa mềm.
      \item Kiểm tra \texttt{IsAvailable == true}.
      \item Kiểm tra slot chưa có appointment \texttt{Pending} hoặc \texttt{Confirmed}.
      \item Khi đặt lịch thành công, set \texttt{IsAvailable = false}.
    \end{itemize}
  \end{block}

  \begin{lstlisting}[language=C]
HasIndex(a => a.AvailabilitySlotId)
  .IsUnique()
  .HasFilter("[Status] IN ('Pending', 'Confirmed')");
  \end{lstlisting}
\end{frame}

\section{Kiến trúc và database}

\begin{frame}{Kiến trúc tổng thể}
  \centering
  \begin{tikzpicture}[
    node distance=1.15cm,
    box/.style={rectangle, rounded corners, draw=Accent, fill=Accent!8, text width=7.2cm, align=center, minimum height=0.85cm},
    side/.style={rectangle, rounded corners, draw=Primary, fill=Primary!8, text width=3.6cm, align=center, minimum height=0.85cm},
    arrow/.style={-Latex, thick, draw=Dark}
  ]
    \node[box] (front) {React 18 + TypeScript + Vite\\Frontend};
    \node[box, below=of front] (api) {ASP.NET Core 8 Web API\\Controllers - Services - DTOs};
    \node[box, below=of api] (db) {SQL Server on MonsterASP};
    \node[side, right=1.25cm of api] (storage) {Supabase Storage\\Avatar/File upload};
    \draw[arrow] (front) -- node[right]{HTTP / JSON / JWT} (api);
    \draw[arrow] (api) -- node[right]{EF Core 8} (db);
    \draw[arrow] (api) -- (storage);
  \end{tikzpicture}
\end{frame}

\begin{frame}{Công nghệ sử dụng}
  \begin{columns}[T]
    \begin{column}{0.32\textwidth}
      \begin{block}{Frontend}
        \begin{itemize}
          \item React 18
          \item TypeScript
          \item Vite
          \item React Router
          \item CSS custom
        \end{itemize}
      \end{block}
    \end{column}
    \begin{column}{0.34\textwidth}
      \begin{block}{Backend}
        \begin{itemize}
          \item ASP.NET Core 8
          \item EF Core 8
          \item SQL Server provider
          \item JWT Bearer
          \item BCrypt
          \item Swagger/OpenAPI
        \end{itemize}
      \end{block}
    \end{column}
    \begin{column}{0.32\textwidth}
      \begin{block}{Infrastructure}
        \begin{itemize}
          \item GitHub Actions
          \item MonsterASP
          \item Remote SQL Server
          \item Supabase Storage
          \item FTP deployment
        \end{itemize}
      \end{block}
    \end{column}
  \end{columns}
\end{frame}

\begin{frame}{Thiết kế database}
  \begin{tabular}{p{0.28\textwidth} p{0.62\textwidth}}
    \toprule
    \textbf{Entity} & \textbf{Vai trò} \\
    \midrule
    \texttt{Users} & Tài khoản, mật khẩu hash, email, avatar, trạng thái hoạt động \\
    \texttt{Roles} & Vai trò: Admin, Student, Lecturer \\
    \texttt{Students} & Hồ sơ sinh viên, liên kết 1-1 với user \\
    \texttt{Lecturers} & Hồ sơ giảng viên, liên kết 1-1 với user \\
    \texttt{AvailabilitySlots} & Khung giờ rảnh của giảng viên \\
    \texttt{Appointments} & Lịch hẹn giữa sinh viên, giảng viên và slot \\
    \texttt{Notifications} & Thông báo nội bộ cho người dùng \\
    \bottomrule
  \end{tabular}
\end{frame}

\begin{frame}{Quan hệ dữ liệu chính}
  \begin{columns}[T]
    \begin{column}{0.52\textwidth}
      \begin{itemize}
        \item Role 1 - N User
        \item User 1 - 1 Student
        \item User 1 - 1 Lecturer
        \item Lecturer 1 - N AvailabilitySlot
        \item Student 1 - N Appointment
        \item Lecturer 1 - N Appointment
        \item User 1 - N Notification
      \end{itemize}
    \end{column}
    \begin{column}{0.46\textwidth}
      \begin{block}{Ràng buộc quan trọng}
        \begin{itemize}
          \item \texttt{AccountName} unique
          \item \texttt{EmailAddress} unique
          \item \texttt{StudentCode} unique
          \item \texttt{LecturerCode} unique
          \item Slot active không được đặt trùng
        \end{itemize}
      \end{block}
    \end{column}
  \end{columns}
\end{frame}

\section{Implementation}

\begin{frame}{Backend design}
  \begin{block}{Tổ chức trách nhiệm}
    \begin{description}
      \item[Controllers] Nhận HTTP request, route và role authorization.
      \item[Services] Xử lý nghiệp vụ và validation.
      \item[Repositories] Truy xuất dữ liệu bằng EF Core.
      \item[DTOs] Chuẩn hóa request/response model.
      \item[Models] Entity map với SQL Server.
      \item[Migrations] Quản lý thay đổi schema database.
    \end{description}
  \end{block}

  \begin{alertblock}{Lợi ích}
    Controller gọn hơn, service dễ test hơn, DTO tránh expose trực tiếp entity database.
  \end{alertblock}
\end{frame}

\begin{frame}{REST API tiêu biểu}
  \scriptsize
  \begin{tabular}{p{0.12\textwidth} p{0.35\textwidth} p{0.16\textwidth} p{0.28\textwidth}}
    \toprule
    \textbf{Method} & \textbf{Endpoint} & \textbf{Role} & \textbf{Mục đích} \\
    \midrule
    POST & \texttt{/api/auth/register-student} & Public & Đăng ký sinh viên \\
    POST & \texttt{/api/auth/register-lecturer} & Public & Đăng ký giảng viên \\
    POST & \texttt{/api/auth/login} & Public & Đăng nhập và nhận JWT \\
    GET & \texttt{/api/lecturers} & Public & Xem danh sách giảng viên \\
    POST & \texttt{/api/availabilityslots} & Lecturer & Tạo slot rảnh \\
    POST & \texttt{/api/appointments} & Student & Đặt lịch hẹn \\
    PUT & \texttt{/api/appointments/\{id\}/status} & Lecturer & Cập nhật trạng thái \\
    GET & \texttt{/api/notifications} & Authorized & Xem thông báo \\
    GET & \texttt{/api/admin/users} & Admin & Quản lý user \\
    \bottomrule
  \end{tabular}
\end{frame}

\begin{frame}{Authentication và authorization}
  \centering
  \begin{tikzpicture}[
    node distance=1cm,
    box/.style={rectangle, rounded corners, draw=Accent, fill=Accent!8, text width=7.2cm, align=center, minimum height=0.75cm},
    arrow/.style={-Latex, thick, draw=Primary}
  ]
    \node[box] (a) {User login};
    \node[box, below=of a] (b) {Backend verify password bằng BCrypt};
    \node[box, below=of b] (c) {Backend tạo JWT token};
    \node[box, below=of c] (d) {Frontend gửi Authorization: Bearer token};
    \node[box, below=of d] (e) {Backend kiểm tra token và role trước khi xử lý API};
    \draw[arrow] (a) -- (b);
    \draw[arrow] (b) -- (c);
    \draw[arrow] (c) -- (d);
    \draw[arrow] (d) -- (e);
  \end{tikzpicture}
\end{frame}

\begin{frame}{Frontend design}
  Frontend được xây dựng bằng React + TypeScript, chia theo page và component.

  \begin{columns}[T]
    \begin{column}{0.5\textwidth}
      \begin{block}{Nhóm màn hình}
        \begin{itemize}
          \item Auth: login, register
          \item Student: profile, lecturer list/detail, appointments
          \item Lecturer: slots, appointment requests
          \item Admin: users, roles
          \item Common: dashboard, notification, protected route
        \end{itemize}
      \end{block}
    \end{column}
    \begin{column}{0.5\textwidth}
      \begin{block}{Điểm triển khai}
        \begin{itemize}
          \item API client tách logic gọi backend.
          \item Protected route bảo vệ màn hình.
          \item TypeScript type giảm lỗi dữ liệu.
          \item Component dùng lại cho shell, badge, route.
        \end{itemize}
      \end{block}
    \end{column}
  \end{columns}
\end{frame}

\begin{frame}{Supabase Storage}
  Project dùng Supabase Storage cho upload avatar/file qua backend.

  \begin{tikzpicture}[
    node distance=1.05cm,
    box/.style={rectangle, rounded corners, draw=Primary, fill=Primary!8, text width=8cm, align=center, minimum height=0.72cm},
    arrow/.style={-Latex, thick, draw=Accent}
  ]
    \node[box] (a) {Frontend gửi multipart file};
    \node[box, below=of a] (b) {ASP.NET Core API nhận \texttt{IFormFile}};
    \node[box, below=of b] (c) {\texttt{SupabaseStorageService} upload qua REST API};
    \node[box, below=of c] (d) {Backend trả về public URL};
    \node[box, below=of d] (e) {URL được lưu hoặc dùng trong profile};
    \draw[arrow] (a) -- (b);
    \draw[arrow] (b) -- (c);
    \draw[arrow] (c) -- (d);
    \draw[arrow] (d) -- (e);
  \end{tikzpicture}

  \vspace{0.25cm}
  \small Giới hạn file size 10 MB và sinh tên file bằng \texttt{Guid} để tránh trùng tên.
\end{frame}

\section{Deployment và đánh giá}

\begin{frame}{CI/CD và deployment}
  Backend deploy tự động qua GitHub Actions.

  \begin{enumerate}
    \item Push vào branch \texttt{main}.
    \item Setup .NET 8 SDK.
    \item Restore dependencies.
    \item Install \texttt{dotnet-ef}.
    \item Test TCP connection tới SQL Server.
    \item Apply EF Core migrations.
    \item Publish ASP.NET Core app.
    \item Inject production connection string.
    \item FTP deploy lên MonsterASP.
  \end{enumerate}

  \begin{block}{Production config}
    Connection string được lưu trong GitHub Secret \texttt{DB\_CONNECTION\_STRING}.
  \end{block}
\end{frame}

\begin{frame}{Khó khăn khi triển khai}
  \begin{block}{SQL Server remote access}
    GitHub Actions chạy migration từ runner bên ngoài, nên SQL Server MonsterASP phải bật remote access port \texttt{1433}.
  \end{block}

  \begin{block}{Secret management}
    Production connection string không nên hard-code trong source code. CI/CD sử dụng GitHub Secret.
  \end{block}

  \begin{block}{Nested Git repository}
    Thư mục \texttt{AcademicAppointmentDeploy} từng bị Git nhận như submodule thiếu \texttt{.gitmodules}, gây warning ở post-job cleanup.
  \end{block}
\end{frame}

\begin{frame}{Kết quả đạt được}
  \begin{itemize}
    \item Xây dựng được workflow đặt lịch giữa sinh viên và giảng viên.
    \item Có đăng ký, đăng nhập, JWT và phân quyền theo role.
    \item Có quản lý slot, appointment, notification và admin user.
    \item Database có ràng buộc chống trùng lịch.
    \item Backend có cấu trúc controller/service/repository rõ ràng.
    \item Frontend tách page, component và API client.
    \item Có pipeline deploy backend lên MonsterASP.
  \end{itemize}
\end{frame}

\begin{frame}{Hạn chế hiện tại}
  \begin{itemize}
    \item Notification hiện là dạng trong hệ thống, chưa real-time.
    \item Chưa có email reminder hoặc calendar integration.
    \item CI/CD chưa chạy toàn bộ unit test trước deploy.
    \item FTP deployment vẫn phụ thuộc hosting truyền thống.
    \item Secret trong local config cần được quản lý kỹ khi public repository.
  \end{itemize}
\end{frame}

\begin{frame}{Hướng phát triển}
  \begin{enumerate}
    \item Tích hợp email reminder khi lịch hẹn được xác nhận.
    \item Đồng bộ Google Calendar hoặc Outlook Calendar.
    \item Real-time notification bằng SignalR.
    \item Bộ lọc nâng cao cho lecturer theo khoa, chuyên ngành, slot rảnh.
    \item Audit log cho thao tác admin và thay đổi trạng thái appointment.
    \item Chạy automated tests trong GitHub Actions trước khi deploy.
    \item Docker hóa backend để deployment ổn định hơn.
  \end{enumerate}
\end{frame}

\section{Demo}

\begin{frame}{Demo flow đề xuất}
  \begin{enumerate}
    \item Login với Student.
    \item Xem danh sách giảng viên.
    \item Chọn giảng viên và slot còn trống.
    \item Tạo appointment.
    \item Login với Lecturer.
    \item Xem request mới.
    \item Confirm hoặc Reject appointment.
    \item Login lại Student.
    \item Kiểm tra status và notification.
    \item Login Admin để xem quản lý user/role.
  \end{enumerate}

  \begin{alertblock}{Trọng tâm demo}
    Student đặt lịch $\rightarrow$ Lecturer xử lý $\rightarrow$ Student nhận trạng thái.
  \end{alertblock}
\end{frame}

\begin{frame}{Kết luận}
  Academic Appointment System giải quyết bài toán đặt lịch tư vấn học thuật bằng một hệ thống web có phân quyền, workflow rõ trạng thái, database constraint và deployment thực tế.

  \begin{block}{Core value}
    Giảm thao tác thủ công, tránh trùng lịch và tập trung hóa quá trình quản lý appointment.
  \end{block}
\end{frame}

\begin{frame}
  \centering
  \Huge Cảm ơn thầy cô đã lắng nghe

  \vspace{0.8cm}
  \Large Q\&A

  \vspace{1cm}
  \normalsize
  Academic Appointment System\\
  ASP.NET Core Web API - React TypeScript - SQL Server
\end{frame}

\end{document}
